%%% The main file. It contains definitions of basic parameters and includes all other parts.

% Meta-data of your thesis (please edit)
\input metadata.tex

% Generate metadata in XMP format for use by the pdfx package
\input xmp.tex

%% Settings for single-side (simplex) printing
% Margins: left 40mm, right 25mm, top and bottom 25mm
% (but beware, LaTeX adds 1in implicitly)
\documentclass[12pt,a4paper]{report}
\setlength\textwidth{145mm}
\setlength\textheight{247mm}
\setlength\oddsidemargin{15mm}
\setlength\evensidemargin{15mm}
\setlength\topmargin{0mm}
\setlength\headsep{0mm}
\setlength\headheight{0mm}
% \openright makes the following text appear on a right-hand page
\let\openright=\clearpage

%% Settings for two-sided (duplex) printing
% \documentclass[12pt,a4paper,twoside,openright]{report}
% \setlength\textwidth{145mm}
% \setlength\textheight{247mm}
% \setlength\oddsidemargin{14.2mm}
% \setlength\evensidemargin{0mm}
% \setlength\topmargin{0mm}
% \setlength\headsep{0mm}
% \setlength\headheight{0mm}
% \let\openright=\cleardoublepage

%% If the thesis has no printed version, symmetric margins look better
% \documentclass[12pt,a4paper]{report}
% \setlength\textwidth{145mm}
% \setlength\textheight{247mm}
% \setlength\oddsidemargin{10mm}
% \setlength\evensidemargin{10mm}
% \setlength\topmargin{0mm}
% \setlength\headsep{0mm}
% \setlength\headheight{0mm}
% \let\openright=\clearpage

%% Generate PDF/A-2u
\usepackage[a-2u]{pdfx}

%% Prefer Latin Modern fonts
\usepackage{lmodern}
% If we are not using LuaTeX, we need to set up character encoding:
\usepackage{iftex}
\ifpdftex
\usepackage[utf8]{inputenc}
\usepackage[T1]{fontenc}
\usepackage{textcomp}
\fi

%% Further useful packages (included in most LaTeX distributions)
\usepackage{amsmath}        % extensions for typesetting of math
\usepackage{amsfonts}       % math fonts
\usepackage{amsthm}         % theorems, definitions, etc.
\usepackage{bm}             % boldface symbols (\bm)
\usepackage{booktabs}       % improved horizontal lines in tables
\usepackage{caption}        % custom captions of floating objects
\usepackage{dcolumn}        % improved alignment of table columns
\usepackage{floatrow}       % custom float environments
\usepackage{graphicx}       % embedding of pictures
\usepackage{indentfirst}    % indent the first paragraph of a chapter
\usepackage[nopatch=item]{microtype}   % micro-typographic refinement
\usepackage{paralist}       % improved enumerate and itemize
\usepackage[nottoc]{tocbibind} % makes sure that bibliography and the lists
			    % of figures/tables are included in the table
			    % of contents
\usepackage{xcolor}         % typesetting in color

% The hyperref package for clickable links in PDF and also for storing
% metadata to PDF (including the table of contents).
% Most settings are pre-set by the pdfx package.
\hypersetup{unicode}
\hypersetup{breaklinks=true}

% Packages for computer science theses
\usepackage{algpseudocode}  % part of algorithmicx package
\usepackage{algorithm}
\usepackage{fancyvrb}       % improved verbatim environment
\usepackage{listings}       % pretty-printer of source code

% You might want to use cleveref for references
% \usepackage{cleveref}

% Set up formatting of bibliography (references to literature)
% Details can be adjusted in macros.tex.
%
% BEWARE: Different fields of research and different university departments
% have their own customs regarding bibliography. Consult the bibliography
% format with your supervisor.
%
% The basic format according to the ISO 690 standard with numbered references
\usepackage[natbib,style=iso-numeric,sorting=none]{biblatex}
% ISO 690 with alphanumeric references (abbreviations of authors' names)
%\usepackage[natbib,style=iso-alphabetic]{biblatex}
% ISO 690 with references Author (year)
%\usepackage[natbib,style=iso-authoryear]{biblatex}
%
% Some fields of research prefer a simple format with numbered references
% (sorting=none tells that bibliography should be listed in citation order)
%\usepackage[natbib,style=numeric,sorting=none]{biblatex}
% Numbered references, but [1,2,3,4,5] is compressed to [1-5]
%\usepackage[natbib,style=numeric-comp,sorting=none]{biblatex}
% A simple format with alphanumeric references:
%\usepackage[natbib,style=alphabetic]{biblatex}

% Load the file with bibliography entries
\addbibresource{bibliography.bib}

% Our definitions of macros (see description inside)
\input macros.tex

%%% Title page and various mandatory informational pages
\begin{document}
%%% Title page of the thesis and other mandatory pages

%%% Inscriptions at the opening page of the hard cover

% We usually do not typeset the hard cover, but if you want to do it, change \iffalse to \iftrue
\iffalse

\pagestyle{empty}
\hypersetup{pageanchor=false}
\begin{center}

\large
Charles University

\medskip

Faculty of Mathematics and Physics

\vfill

{\huge\bf\ThesisTypeTitle}

\vfill

{\huge\bf\ThesisTitle\par}

\vfill
\vfill

\hbox to \hsize{\YearSubmitted\hfil \ThesisAuthor}

\end{center}

\newpage\openright
\setcounter{page}{1}

\fi

%%% Title page of the thesis

\pagestyle{empty}
\hypersetup{pageanchor=false}
\begin{center}

\centerline{\mbox{\includegraphics[width=166mm]{img/logo-en.pdf}}}

\vspace{-8mm}
\vfill

{\bf\Large\ThesisTypeTitle}

\vfill

{\LARGE\ThesisAuthor}

\vspace{15mm}

{\LARGE\bfseries\ThesisTitle\par}

\vfill

\Department

\vfill

{
\centerline{\vbox{\halign{\hbox to 0.45\hsize{\hfil #}&\hskip 0.5em\parbox[t]{0.45\hsize}{\raggedright #}\cr
Supervisor of the \ThesisTypeName{} thesis:&\Supervisor \cr
\ifx\ThesisType\TypeRig\else
\noalign{\vspace{2mm}}
Study programme:&\StudyProgramme \cr
\fi
}}}}

\vfill

Prague \YearSubmitted

\end{center}

\newpage

%%% A page with a solemn declaration to the thesis

\openright
\hypersetup{pageanchor=true}
\vglue 0pt plus 1fill

\noindent
I declare that I carried out this \ThesisTypeName{} thesis on my own, and only with the cited
sources, literature and other professional sources.
I understand that my work relates to the rights and obligations under the Act No.~121/2000 Sb.,
the Copyright Act, as amended, in particular the fact that the Charles
University has the right to conclude a license agreement on the use of this
work as a school work pursuant to Section 60 subsection 1 of the Copyright~Act.

\vspace{10mm}

\hbox{\hbox to 0.5\hsize{%
In \hbox to 6em{\dotfill} date \hbox to 6em{\dotfill}
\hss}\hbox to 0.5\hsize{\dotfill\quad}}
\smallskip
\hbox{\hbox to 0.5\hsize{}\hbox to 0.5\hsize{\hfil Author's signature\hfil}}

\vspace{20mm}
\newpage

%%% Dedication

\openright

\noindent
\Dedication

\newpage

%%% Mandatory information page of the thesis

\openright
{\InfoPageFont

\vtop to 0.5\vsize{
\setlength\parindent{0mm}
\setlength\parskip{5mm}

Title:
\ThesisTitle

Author:
\ThesisAuthor

\DeptType:
\Department

Supervisor:
\Supervisor, \SupervisorsDepartment

Abstract:
\Abstract

Keywords:
{\def\sep{\unskip, }\ThesisKeywords}

\vfil
}

% In Czech study programmes, it is mandatory to include Czech meta-data:

\ifx\StudyLanguage\LangCS

\vtop to 0.49\vsize{
\setlength\parindent{0mm}
\setlength\parskip{5mm}

Název práce:
\ThesisTitleCS

Autor:
\ThesisAuthor

\DeptTypeCS:
\DepartmentCS

Vedoucí diplomové práce:
\Supervisor, \SupervisorsDepartmentCS

Abstrakt:
\AbstractCS

Klíčová slova:
{\def\sep{\unskip, }\ThesisKeywordsCS}

\vfil
}

\fi

}

\newpage

%%% Further pages will be numbered
\pagestyle{plain}


%%% A page with automatically generated table of contents of the thesis

\tableofcontents

%%% Each chapter is kept in a separate file
\chapter*{Introduction}
\addcontentsline{toc}{chapter}{Introduction}


\section*{Motivation}

Domain modeling is fundamental both in data and software engineering. Domain modeling techniques, such as conceptual modeling techniques, where a domain model is expressed as an ER or UML class diagram, or ontology engineering techniques, which define a domain model as a domain ontology, are well established both in academia and industry \cite{Verdonck2018}. However, creating these domain models manually takes a non-trivial amount of time and requires trained experts. Therefore, various modeling automation techniques have been proposed \cite{Sonbol2022}. Recently, researchers began to evaluate different approaches to using large language models (LLMs) as possible automation solutions \cite{Chen2023,Saeedizade2024}.


\section*{Solution approach discussion}

Currently, fully automated domain modeling with LLMs works only for simple domain descriptions with a small number of domain elements \cite{Camara2023}. This is true even when advanced prompting techniques are applied \cite{Chen2023,Saeedizade2024}. To mitigate this issue, the user can be introduced into the domain modeling process \cite{Camara2023} and the LLM can be used as an assistant. For example, \citet{Camara2023} experimented with domain modeling in a conversation dialogue fashion with the ChatGPT-3.5. First, the LLM generated some preliminary domain model and then the authors created follow-up instructions to edit their domain models. However, as the LLM executed directly each authors' instruction, this frequently led to back-and-forth communication as the authors tried correcting the mistakes made by the LLM in the previous steps. Additionally, there was no option to manually model any parts of the domain model to help the LLM model the parts it struggled with. For a detailed description of related work see chapter \ref{chap:related_work}.


\section*{Using LLM as an assistant}
\label{section:llm_as_an_assistant}

Recently, a lot of applications started using the LLM as an assistant. Here is a list of some of those applications with a brief description of the assistant usage.

\begin{itemize}
\item \textit{Grammarly}\footnote{\url{https://app.grammarly.com/}}: assistant for writing
\item \textit{GitHub Copilot}\footnote{\url{https://github.com/features/copilot}}: assistant for programming
\item \textit{Notion AI}\footnote{\url{https://www.notion.so/help/guides/category/ai}}: assistant for taking notes
\end{itemize}

These applications have the following features in common. First, the LLM usually provides only suggestions to the users. These suggestions usually can be accepted or rejected or in some cases, a new suggestion can be re-generated. The reason for this behaviour is that the LLMs still make a lot of mistakes so direct application of their output is usually not wanted. Typically, these applications explicitly warn their users that the generated output can provide false, biased and outdated information. Furthermore, if these mistakes are not removed by the users these mistakes can accumulate over time and can lead to even worse output performance. Second, the suggestions are usually in a form such that the users can quickly decide if they want to use them. For example, in \textit{Grammarly} the spell-checking feature represents each suggestion as an underlined text that contains a reason why it was underlined and a suggested correction. In \textit{GitHub Copilot} especially in the hands of an experienced developer, it is usually easy to decide if a suggested piece of code or a suggested code documentation is appropriate based on the surrounding context. And finally, in \textit{Notion AI} suggested notes are usually easy to evaluate as they are typically clearly structured.


\section*{Areas with lack of research}

Based on our exploration, the following areas have a lack of research and can have a significant impact on making the domain modeling process faster and more accessible:

\begin{itemize}
\item manual domain modeling extended with the LLM as an assistant that only provides suggestions
\item using techniques for filtering domain descriptions such as retrieval-augmented generation to improve the performance of LLMs for the domain modeling tasks
\item using real-life domain descriptions for evaluating domain modeling tasks with LLMs
\item creating a tool for demonstrating the viability of the selected approach
\end{itemize}


\section*{Goals}

The domain modeling process typically follows a systematic approach. First, the domain experts create an unstructured domain description, and second, the modeling experts solely based on this domain description create an exact, structured, and formalized domain model.

This thesis aims to help with the second step and address the mentioned areas with lack of research by creating a tool for domain modeling where the LLM assistant cooperates with a human modeling expert. The modeling expert provides an unstructured domain description as input into the domain modeling process and the assistant generates suggestions solely based on this domain description.

To improve the performance of the LLM, we divide the domain modeling process into a set of simpler steps and we employ the retrieval-augmented generation technique for providing only the relevant parts of the domain description based on the given task to the LLM. This assistant always generates output in the form of suggestions as the LLMs still frequently make mistakes. To help the users decide if they want to apply the generated suggestions, we highlight each suggested domain element in the given domain description. Finally, we experimentally evaluate our approach with a real-life domain descriptions.

This LLM assistant is mainly intended to help the users create their domain models faster and to help them better document their domain models. This assistant should be:

\begin{enumerate}
\item intuitive to use
\item suggesting mostly appropriate classes, attributes, and associations
\item suggesting mostly appropriate descriptions of classes, attributes, and associations
\item able to summarize a selected part of a domain model
\item able to highlight in the given domain description the parts that the user already modeled
\item able to generate corresponding suggestions in a few seconds
\end{enumerate}


\section*{Outline}

Chapter \ref{chap:background} provides an overview of the basic technologies that we will be using.
Chapter \ref{chap:domain_modeling_process_analysis} introduces the minimalist but necessary formalization of the domain modeling process and its steps that can be automated.
Chapter \ref{chap:framework} describes a theoretical generic framework for building LLM-based domain modeling assistants.
Chapter \ref{chap:framework_configuration} shows some possible configurations of the framework.
Chapter \ref{chap:protype_software_application} implements the framework with the proposed configurations.
Chapter \ref{chap:evaluation} presents the experimental evaluation of our framework implementation.
Chapter \ref{chap:development_documentation} contains the development documentation and chapter \ref{chap:user_documentation} contains the user documentation. Chapter \ref{chap:related_work} gives an overview of the related work and chapter \ref{chap:conclusion} concludes our work and suggests possible future work.
\chapter{Data}

%An~example citation: \cite{Andel07}



\chapter{Methodology}


\section{Prompting methods}

\subsection{What are they}

Něco jako: "In the field of LLMs, a prompt is an input that guides the model’s response
 generation. A prompting technique entails the strategic formulation of these
 prompts to maximize the efficacy of LLMs. It involves the deliberate structuring
 and phrasing of prompts to align with the model’s training and capabilities." (zdroj: Navigating Ontology Development with LLMs)

\subsection{Prompt structure}

We created prompts in form of templates where before sending this prompt to the LLM each symbol is replaced by a given argument.


\subsection{Prompting techniques}
\begin{itemize}
\item subsections: CoT, ToT, few-shot prompting \\
\item CoT can be further divided into more approaches
\end{itemize}



\section{RAG methods}
- TODO: brief explanation of what RAG is, how it works and why it works \\
- RAG approaches: semantic vs. syntactic (more in GAO, Yunfan, et al. Retrieval-augmented generation for large language models: A survey.) \\
\chapter{Evaluation}

\section{Data}
%An~example citation: \cite{Andel07}

In our application we work with domain descriptions to generate suggestions. \\

Data, se kterými pracujeme:
\begin{itemize}
\item popisy domén
\item anotované popisy domén
\end{itemize}


TODO: Možná bude stačit dát k evaluaci sekci o datech, protože toho ve výsledku o těch datech neřeknu moc:
\begin{itemize}
\item jak data vypadají, tedy jednotlivé styly (educational, atd.)
\item možná zmínit náš primární use-case: analytik, který nejprve popis domény příslušným způsobem upraví
\item možná jak byla vytvořena (především ručně, nevygeneroval to LLM)
\item kolik mají tříd, atributů, asociací
\item jeden popis domény používáme pro N-shot prompting
\item zdůvodnit, k čemu máme anotované popisy domén
\end{itemize}



\chapter{Future work}

- dynamic domain description \\

- adding attribute as a relationship (and vice versa) and letting LLM suggest a new appropriate name for this attribute/relationship \\

- translator of user's plain text instructions into changes inside a conceptual model \\

- in domain description option to highlight the part that the assistant should solely focus on \\

- conceptual model updating: assistant that based on some instruction and given conceptual model creates an updated conceptual model (more info: consultation 23) \\

- suggestions generation based on chosen ontology design patterns (more info: consultation 23) \\



\chapter*{Conclusion}
\addcontentsline{toc}{chapter}{Conclusion}


%%% Bibliography
%%% Bibliography (literature used as a source)
%%%
%%% We employ biblatex to construct the bibliography. It processes
%%% citations in the text (e.g., the \cite{...} macro) and looks up
%%% relevant entries in the bibliography.bib file.
%%%
%%% See also biblatex settings in thesis.tex.

%%% Generate the bibliography. Beware that if you cited no works,
%%% the empty list will be omitted completely.

% We let bibliography items stick out of the right margin a little
\def\bibfont{\hfuzz=2pt}

\printbibliography[heading=bibintoc]

%%% If case you prefer to write the bibliography manually (without biblatex),
%%% you can use the following. Please follow the ISO 690 standard and
%%% citation conventions of your field of research.

% \begin{thebibliography}{99}
%
% \bibitem{lamport94}
%   {\sc Lamport,} Leslie.
%   \emph{\LaTeX: A Document Preparation System}.
%   2nd edition.
%   Massachusetts: Addison Wesley, 1994.
%   ISBN 0-201-52983-1.
%
% \end{thebibliography}


%%% Figures used in the thesis (consider if this is needed)
\listoffigures

%%% Tables used in the thesis (consider if this is needed)
%%% In mathematical theses, it could be better to move the list of tables to the beginning of the thesis.
\listoftables

%%% Abbreviations used in the thesis, if any, including their explanation
%%% In mathematical theses, it could be better to move the list of abbreviations to the beginning of the thesis.
%\chapwithtoc{List of Abbreviations}

%%% Doctoral theses must contain a list of author's publications
\ifx\ThesisType\TypePhD
\chapwithtoc{List of Publications}
\fi

%%% Attachments to the thesis, if any. Each attachment must be referred to
%%% at least once from the text of the thesis. Attachments are numbered.
%%%
%%% The printed version should preferably contain attachments, which can be
%%% read (additional tables and charts, supplementary text, examples of
%%% program output, etc.). The electronic version is more suited for attachments
%%% which will likely be used in an electronic form rather than read (program
%%% source code, data files, interactive charts, etc.). Electronic attachments
%%% should be uploaded to SIS. Allowed file formats are specified in provision
%%% of the rector no. 72/2017. Exceptions can be approved by faculty's coordinator.
\appendix
%\chapter{Attachments}

%\section{First Attachment}

\end{document}
